% GENERATED from WHITEPAPER.md by scripts/whitepaper/md2tex.sh. Do not edit by hand.
% Options for packages loaded elsewhere
\PassOptionsToPackage{unicode}{hyperref}
\PassOptionsToPackage{hyphens}{url}
\PassOptionsToPackage{dvipsnames,svgnames,x11names}{xcolor}
\documentclass[
  11pt,
  a4paper]{article}
\usepackage{xcolor}
\usepackage{amsmath,amssymb}
\setcounter{secnumdepth}{-\maxdimen} % remove section numbering
\usepackage{iftex}
\ifPDFTeX
  \usepackage[T1]{fontenc}
  \usepackage[utf8]{inputenc}
  \usepackage{textcomp} % provide euro and other symbols
\else % if luatex or xetex
  \usepackage{unicode-math} % this also loads fontspec
  \defaultfontfeatures{Scale=MatchLowercase}
  \defaultfontfeatures[\rmfamily]{Ligatures=TeX,Scale=1}
\fi
\usepackage[]{lmodern}
\ifPDFTeX\else
  % xetex/luatex font selection
\fi
% Use upquote if available, for straight quotes in verbatim environments
\IfFileExists{upquote.sty}{\usepackage{upquote}}{}
\IfFileExists{microtype.sty}{% use microtype if available
  \usepackage[]{microtype}
  \UseMicrotypeSet[protrusion]{basicmath} % disable protrusion for tt fonts
}{}
\makeatletter
\@ifundefined{KOMAClassName}{% if non-KOMA class
  \IfFileExists{parskip.sty}{%
    \usepackage{parskip}
  }{% else
    \setlength{\parindent}{0pt}
    \setlength{\parskip}{6pt plus 2pt minus 1pt}}
}{% if KOMA class
  \KOMAoptions{parskip=half}}
\makeatother
\usepackage{listings}
\newcommand{\passthrough}[1]{#1}
\lstset{defaultdialect=[5.3]Lua}
\lstset{defaultdialect=[x86masm]Assembler}
\usepackage{longtable,booktabs,array}
\usepackage{caption}
\captionsetup[table]{skip=6pt}
\newcounter{none} % for unnumbered tables
\usepackage{calc} % for calculating minipage widths
% Correct order of tables after \paragraph or \subparagraph
\usepackage{etoolbox}
\makeatletter
\patchcmd\longtable{\par}{\if@noskipsec\mbox{}\fi\par}{}{}
\makeatother
\setlength{\emergencystretch}{3em} % prevent overfull lines
\providecommand{\tightlist}{%
  \setlength{\itemsep}{0pt}\setlength{\parskip}{0pt}}
% Preamble additions for WHITEPAPER.tex, injected by scripts/whitepaper/md2tex.sh.
\usepackage{microtype}
\usepackage[margin=1in]{geometry}
\usepackage{xcolor}
\usepackage{textcomp}


\lstdefinelanguage{Rust}{
  keywords={fn, let, mut, struct, enum, impl, for, in, if, else, return, match, use,
           pub, crate, mod, self, Some, None, Ok, Err, as, where, dyn, trait, type,
           const, static, ref, move, async, await, break, continue, loop, while,
           true, false},
  keywords=[2]{Vec, HashMap, HashSet, Option, Result, String, Box, Uuid, DateTime,
           Utc, usize, u64, u32, u8, f32, f64, bool, str},
  morecomment=[l]{//},
  morecomment=[s]{/*}{*/},
  morestring=[b]",
  sensitive=true,
}

\lstset{
  basicstyle=\ttfamily\footnotesize,
  keywordstyle=\color{blue!60!black}\bfseries,
  keywordstyle=[2]\color{teal!70!black},
  commentstyle=\color{green!40!black}\itshape,
  stringstyle=\color{red!55!black},
  showstringspaces=false,
  breaklines=true,
  breakatwhitespace=false,
  frame=single,
  framerule=0.4pt,
  rulecolor=\color{gray!50},
  backgroundcolor=\color{gray!5},
  tabsize=2,
  columns=fixed,
  basewidth=0.52em,
  keepspaces=true,
  upquote=true,
  aboveskip=0.8em,
  belowskip=0.8em,
  xleftmargin=0.5em,
  xrightmargin=0.5em,
}

% Long tables: slightly tighter rows and smaller text.
\AtBeginEnvironment{longtable}{\small}
\renewcommand{\arraystretch}{1.15}
\usepackage{bookmark}
\IfFileExists{xurl.sty}{\usepackage{xurl}}{} % add URL line breaks if available
\urlstyle{same}
% fallback for those not using the hyperref driver hyperxmp:
\makeatletter
\@ifundefined{xmpquote}{\newcommand{\xmpquote}[1]{#1}}{}
\makeatother
\hypersetup{
  colorlinks=true,
  linkcolor={black},
  filecolor={Maroon},
  citecolor={Blue},
  urlcolor={blue},
  pdfcreator={LaTeX via pandoc}}

\title{\textbf{MentisDB}\\[0.3em]\large Durable, Verifiable Memory for Software Agents}
\author{Angel Leon}
\date{master after 0.10.10.55 (unreleased changes included) --- 2026-11-01}

\begin{document}
\maketitle
\begin{abstract}
Agents forget. A coding agent that learned last month that the user
wants tabs, that the build breaks on a certain linker flag, or that a
design was rejected, starts the next session without any of it unless
something durable and searchable carries it forward. MentisDB is that
memory, built as an embedded Rust library and an optional single-binary
daemon. Memory is an append-only log of typed records called
\emph{thoughts}. Each record is linked to its predecessor by a SHA-256
hash, can carry an Ed25519 signature bound to its chain and position,
and can be governed by a write policy recorded inside the chain itself.

Retrieval combines BM25 lexical scoring with document-frequency gating,
dense-vector similarity from a local ONNX model, typed graph expansion,
append-order session cohesion, and optional Reciprocal Rank Fusion.
Nothing is deleted: a superseded thought is replaced in results by its
current successor, and one that a newer thought merely contests is shown
with the contester ranked above it. A \emph{standing context} query
returns what is durably true about the user and the project regardless
of age, next to the usual window of recent thoughts. None of this needs
an external database, a network service or an LLM call.

On the LoCoMo-10P benchmark the default pipeline reaches 73.5\%
Recall@10, and on LongMemEval it reaches 68.0\% Recall@5 (release
0.10.9.54, local MiniLM embeddings). On real project chains, an audit
found no loss of recall with the age of a memory, but it found that
near-duplicate thoughts crowd results and that paraphrased queries
recall far less than verbatim ones. This paper describes the design as
it is implemented, states each default exactly, and lists what is still
weak.
\end{abstract}

\subsection{1. The Problem and the Requirements}

An agent's context window is working memory. It is erased at the end of
each session and compacted during long ones. What the agent needs on top
of it is the equivalent of long-term memory: a store it writes to as it
works and reads from when it starts or when a question touches the past.
We set eight requirements for that store.

\begin{enumerate}
\def\labelenumi{\arabic{enumi}.}
\tightlist
\item
  \textbf{Durable.} An acknowledged write survives a process crash. With
  an opt-in setting it also survives power loss.
\item
  \textbf{Recallable regardless of age.} A preference stated three
  months ago must reach the agent as reliably as one stated yesterday. A
  window of recent records does not satisfy this on its own.
\item
  \textbf{Searchable by meaning and by words.} Users paraphrase. Exact
  identifiers, file names and error strings also matter. Both kinds of
  query must work.
\item
  \textbf{Current.} When a fact changes, the new version must win and
  the old one must stop surfacing, without destroying the history that
  explains the change.
\item
  \textbf{Attributable and verifiable.} Every record names the agent
  that wrote it. Edits to history are detectable, and a chain can
  require that writes carry valid signatures.
\item
  \textbf{Local and dependency-free.} The core path runs on a laptop
  with no database server, no cloud API and no per-write LLM cost.
\item
  \textbf{Portable across agents.} Claude Code, Codex, Copilot, Gemini
  and any other MCP client share one memory.
\item
  \textbf{Inspectable.} Every ranking can be explained, and every memory
  can be read, exported and audited by a human.
\end{enumerate}

Common approaches fail at least one requirement. A Markdown notes file
is durable and portable but is not searchable at scale and is loaded
whole or not at all. Provider session memory is not portable and not
inspectable. A plain vector store answers paraphrases but has no notion
of currency, attribution or integrity, and misses exact identifiers.
Systems that extract memories with an LLM on every write add latency,
cost and a second model whose mistakes become permanent.

MentisDB addresses the requirements with three ideas. Memory is an
append-only, hash-chained log, so history is never rewritten and changes
are recorded as new thoughts that point at old ones. Retrieval is a
local, multi-signal ranking whose every component is exposed in the
result. A written protocol, shipped as an agent skill, tells agents when
to read, when to write and what to link.


\subsection{2. Architecture}

MentisDB is a Rust crate (\passthrough{\lstinline!mentisdb!}) and one
binary of the same name. The crate can be embedded directly. The binary
runs a daemon and also serves as the command-line client.

\begin{lstlisting}
  Agents (Claude Code, Codex, Copilot, Gemini, ...)      Humans
        |  MCP (HTTP or stdio)    |  REST       |  CLI      |  browser
        v                         v             v           v
  +---------------------------------------------------------------+
  | mentisdb daemon (one process, binds 127.0.0.1 by default)     |
  |   MCP 9471   REST 9472   HTTPS MCP 9473   HTTPS REST 9474     |
  |   Dashboard 9475 (HTTPS)                                      |
  |   auth: optional bearer tokens (global or chain-scoped)       |
  +-------------------------------+-------------------------------+
                                  |
  +-------------------------------v-------------------------------+
  | MentisDb core (library)                                       |
  |   append path: validate -> access policy -> sign check ->     |
  |                relations/dedup -> hash -> writer thread       |
  |   read path:   candidates (invalidation, as_of, filters) ->   |
  |                BM25 + vectors + graph + cohesion -> rank      |
  |   standing context, recent context, context bundles           |
  |   agent registry, entity types, skill registry, webhooks      |
  +-------------------------------+-------------------------------+
                                  |
  +-------------------------------v-------------------------------+
  | Files under MENTISDB_DIR (default ~/.cloudllm/mentisdb)       |
  |   canonical:  <stem>.tcbin (the hash chain)                   |
  |   sidecars:   agents, entity types, vectors (+WAL, HNSW),     |
  |               implicit-edge overlay   (all rebuildable)       |
  |   global:     chain registry, skills, webhooks, tokens, TLS   |
  +---------------------------------------------------------------+
\end{lstlisting}

The canonical log is the only source of truth. Every other file is
either a small registry written atomically or a derived index that can
be rebuilt from the log.

Ports are set in \passthrough{\lstinline!src/server.rs!}
(\passthrough{\lstinline!MentisDbServerConfig!}) and can be changed or
disabled with \passthrough{\lstinline!MENTISDB\_MCP\_PORT!},
\passthrough{\lstinline!MENTISDB\_REST\_PORT!},
\passthrough{\lstinline!MENTISDB\_HTTPS\_MCP\_PORT!},
\passthrough{\lstinline!MENTISDB\_HTTPS\_REST\_PORT!} and
\passthrough{\lstinline!MENTISDB\_DASHBOARD\_PORT!} (0 disables a
listener). A self-signed TLS certificate is generated on first start
under \passthrough{\lstinline!tls/!}.


\subsection{3. Data Model}

\subsubsection{3.1 Thoughts}

A thought is one memory record. The \passthrough{\lstinline!Thought!}
struct (\passthrough{\lstinline!src/lib.rs!}) has these fields, in
declared order. The order matters because records are encoded
positionally.

{\def\LTcaptype{none} % do not increment counter
\begin{longtable}[]{@{}
  >{\raggedright\arraybackslash}p{(\linewidth - 6\tabcolsep) * \real{0.1881}}
  >{\raggedright\arraybackslash}p{(\linewidth - 6\tabcolsep) * \real{0.2178}}
  >{\raggedright\arraybackslash}p{(\linewidth - 6\tabcolsep) * \real{0.0792}}
  >{\raggedright\arraybackslash}p{(\linewidth - 6\tabcolsep) * \real{0.5149}}@{}}
\toprule\noalign{}
\begin{minipage}[b]{\linewidth}\raggedright
Field
\end{minipage} & \begin{minipage}[b]{\linewidth}\raggedright
Type
\end{minipage} & \begin{minipage}[b]{\linewidth}\raggedright
Set by
\end{minipage} & \begin{minipage}[b]{\linewidth}\raggedright
Meaning
\end{minipage} \\
\midrule\noalign{}
\endhead
\bottomrule\noalign{}
\endlastfoot
\passthrough{\lstinline!schema\_version!} &
\passthrough{\lstinline!u32!} & chain & Record layout version (currently
3) \\
\passthrough{\lstinline!id!} & \passthrough{\lstinline!Uuid!} & chain &
Stable identifier (random v4) \\
\passthrough{\lstinline!index!} & \passthrough{\lstinline!u64!} & chain
& Position in the chain, from 0 \\
\passthrough{\lstinline!timestamp!} &
\passthrough{\lstinline!DateTime<Utc>!} & chain & Commit time \\
\passthrough{\lstinline!session\_id!} &
\passthrough{\lstinline!Option<Uuid>!} & caller & Groups thoughts of one
working session \\
\passthrough{\lstinline!agent\_id!} & \passthrough{\lstinline!String!} &
caller & Author, resolved through the agent registry \\
\passthrough{\lstinline!signing\_key\_id!} &
\passthrough{\lstinline!Option<String>!} & caller & Key that signed this
thought \\
\passthrough{\lstinline!thought\_signature!} &
\passthrough{\lstinline!Option<Vec<u8>>!} & caller & Detached Ed25519
signature (Section 4.2) \\
\passthrough{\lstinline!thought\_type!} &
\passthrough{\lstinline!ThoughtType!} & caller & What kind of memory
this is \\
\passthrough{\lstinline!role!} & \passthrough{\lstinline!ThoughtRole!} &
caller & How the system should treat it \\
\passthrough{\lstinline!content!} & \passthrough{\lstinline!String!} &
caller & The memory text \\
\passthrough{\lstinline!confidence!} &
\passthrough{\lstinline!Option<f32>!} & caller & 0 to 1, clamped \\
\passthrough{\lstinline!importance!} & \passthrough{\lstinline!f32!} &
caller & 0 to 1, clamped; default 0.5 \\
\passthrough{\lstinline!tags!} & \passthrough{\lstinline!Vec<String>!} &
caller & Free-form labels, normalized; also carries the scope \\
\passthrough{\lstinline!concepts!} &
\passthrough{\lstinline!Vec<String>!} & caller & Semantic labels,
normalized \\
\passthrough{\lstinline!refs!} & \passthrough{\lstinline!Vec<u64>!} &
caller & Back-references by index \\
\passthrough{\lstinline!relations!} &
\passthrough{\lstinline!Vec<ThoughtRelation>!} & caller & Typed edges
(Section 3.3) \\
\passthrough{\lstinline!entity\_type!} &
\passthrough{\lstinline!Option<String>!} & caller & Optional domain
label, e.g.~\passthrough{\lstinline!bug\_report!} \\
\passthrough{\lstinline!source\_episode!} &
\passthrough{\lstinline!Option<String>!} & caller & Optional link to an
external source \\
\passthrough{\lstinline!prev\_hash!} & \passthrough{\lstinline!String!}
& chain & Hash of the previous thought, empty for the first \\
\passthrough{\lstinline!hash!} & \passthrough{\lstinline!String!} &
chain & SHA-256 of this record (Section 4.1) \\
\end{longtable}
}

The caller supplies the semantic fields. The chain assigns identity,
position, time and the hash fields, so a client cannot choose where in
history its record lands. For an unsigned thought the chain also
normalizes input: it adds a \passthrough{\lstinline!References!}
relation for each \passthrough{\lstinline!ref!}, fills a missing
\passthrough{\lstinline!valid\_at!} on relations with the commit time,
and may add a \passthrough{\lstinline!Supersedes!} edge when automatic
deduplication is on (Section 6.9). A signed thought is stored exactly as
signed, with none of these additions, so its signature can be checked
again later.

\subsubsection{3.2 Types and Roles}

\passthrough{\lstinline!ThoughtType!} says what kind of memory a thought
is. It has 32 variants. In declared order they are:

\begin{quote}
PreferenceUpdate, UserTrait, RelationshipUpdate, Finding, Insight,
FactLearned, PatternDetected, Hypothesis, Mistake, Correction,
LessonLearned, AssumptionInvalidated, Constraint, Plan, Subgoal,
Decision, StrategyShift, Wonder, Question, Idea, Experiment,
ActionTaken, TaskComplete, Checkpoint, StateSnapshot, Handoff, Summary,
Surprise, Reframe, Goal, LLMExtracted, AccessPolicy.
\end{quote}

\passthrough{\lstinline!LLMExtracted!} marks memories produced by the
optional extraction feature. \passthrough{\lstinline!AccessPolicy!} is
reserved for governance records and is refused on the normal append
path.

\passthrough{\lstinline!ThoughtRole!} says how the system treats a
thought: \passthrough{\lstinline!Memory!},
\passthrough{\lstinline!WorkingMemory!},
\passthrough{\lstinline!Summary!},
\passthrough{\lstinline!Compression!},
\passthrough{\lstinline!Checkpoint!}, \passthrough{\lstinline!Handoff!},
\passthrough{\lstinline!Audit!} or
\passthrough{\lstinline!Retrospective!}. Type and role are independent.
A context-compaction checkpoint is a \passthrough{\lstinline!Summary!}
thought with role \passthrough{\lstinline!Checkpoint!}. A post-mortem
lesson is a \passthrough{\lstinline!LessonLearned!} thought with role
\passthrough{\lstinline!Retrospective!}.

Types drive behaviour, not only labels. Standing context (Section 6.10)
selects thoughts by type.

\subsubsection{3.3 Typed Relations}

A relation is a directed, typed edge from the new thought to an earlier
one, possibly in another chain:

\begin{lstlisting}
struct ThoughtRelation {
    kind: ThoughtRelationKind,
    target_id: Uuid,
    chain_key: Option<String>,        // set for cross-chain edges
    valid_at: Option<DateTime<Utc>>,  // when the edge starts to hold
    invalid_at: Option<DateTime<Utc>>,// when it stops holding
}
\end{lstlisting}

There are twelve kinds: \passthrough{\lstinline!References!},
\passthrough{\lstinline!Summarizes!},
\passthrough{\lstinline!Corrects!},
\passthrough{\lstinline!Invalidates!},
\passthrough{\lstinline!CausedBy!}, \passthrough{\lstinline!Supports!},
\passthrough{\lstinline!Contradicts!},
\passthrough{\lstinline!DerivedFrom!},
\passthrough{\lstinline!ContinuesFrom!},
\passthrough{\lstinline!BranchesFrom!},
\passthrough{\lstinline!RelatedTo!} and
\passthrough{\lstinline!Supersedes!}. Because edges always point
backwards in time, the relation graph of a chain is acyclic by
construction.

Three kinds change what retrieval shows:
\passthrough{\lstinline!Supersedes!}, \passthrough{\lstinline!Corrects!}
and \passthrough{\lstinline!Invalidates!}. A thought that is the target
of one of them is \emph{invalidated} (Section 3.6).

\subsubsection{3.4 Scopes, Sessions and Entities}

A thought belongs to one memory scope: \passthrough{\lstinline!User!}
(the default), \passthrough{\lstinline!Session!} or
\passthrough{\lstinline!Agent!}. The scope is stored as a tag
(\passthrough{\lstinline!scope:user!},
\passthrough{\lstinline!scope:session!},
\passthrough{\lstinline!scope:agent!}), so it is covered by signatures
and filters like any tag. User scope holds facts about the person and
the project. Session and agent scopes hold material private to one
session or one agent, and standing context skips them.

\passthrough{\lstinline!session\_id!} groups the thoughts of one working
session. \passthrough{\lstinline!entity\_type!} lets a chain declare
domain types such as \passthrough{\lstinline!bug\_report!} or
\passthrough{\lstinline!customer!}; the per-chain entity-type registry
records them. Both are filters. Neither changes ranking.

\subsubsection{3.5 Agents}

Thoughts carry only an \passthrough{\lstinline!agent\_id!}. Everything
else about an agent lives in a per-chain registry
(\passthrough{\lstinline!AgentRecord!}): display name, owner,
description, aliases, status (\passthrough{\lstinline!Active!} or
\passthrough{\lstinline!Revoked!}), public keys, first and last seen
index and time, and a thought count. Each public key records its id,
algorithm (Ed25519 is the only one), key bytes, and when it was added
and revoked.

A public key is accepted only if it is a valid, non-weak Ed25519 curve
point and comes with a proof of possession: a signature by that key over
a payload with domain
\passthrough{\lstinline!mentisdb/key-registration/v1!} that names the
chain, agent, key id, key and a timestamp within ten minutes of the
server clock (\passthrough{\lstinline!KeyRegistrationProof::verify!} in
\passthrough{\lstinline!src/access.rs!}). Nobody can register a key they
do not hold.

The registry is a sidecar file, not part of the hash chain. Section 4.3
explains what that means for trust.

\subsubsection{3.6 Versions, Invalidation and Point-in-Time Reads}

Nothing is deleted. When a thought is superseded, corrected or
invalidated, the new thought carries the edge and the old one stays in
the log.

At open, the chain scans all relations and builds the set of invalidated
thought ids together with each target's invalidating edges
(\passthrough{\lstinline!observe\_invalidating\_relations!}). Every
append updates the set. The default read paths hide invalidated
thoughts: filter search, ranked search, context bundles and recent
context. A caller sees them with
\passthrough{\lstinline!include\_invalidated=true!}, which is how an
agent audits what changed.

Hiding alone is not enough, for two reasons. A successor often uses
different words than the thought it replaces, so hiding the old version
can leave the agent with nothing. And agents rarely type their links. On
one real installation, none of 14 \passthrough{\lstinline!Correction!}
and \passthrough{\lstinline!AssumptionInvalidated!} thoughts and none of
about 90 \passthrough{\lstinline!Decision!} thoughts carried a
\passthrough{\lstinline!Supersedes!}, \passthrough{\lstinline!Corrects!}
or \passthrough{\lstinline!Invalidates!} relation; half the corrections
carried no link and the rest only plain references. The invalidated set
stayed empty and the old policy kept winning searches. The lineage index
(\passthrough{\lstinline!src/lineage.rs!}) addresses both problems. It
is derived from the stored chain at open time, so existing chains
benefit without migration.

\begin{itemize}
\tightlist
\item
  \textbf{Superseded.} A thought targeted by a strong edge is
  \passthrough{\lstinline!superseded!}. When it would have ranked,
  ranked search puts its \emph{live successor} in its place: the newest
  unsuperseded end of the chain of supersessions, followed up to 64
  steps. The successor inherits the old thought's score and reports
  which thoughts it stands in for
  (\passthrough{\lstinline!successor\_of!}), so it reaches the agent
  even when its wording does not match the query.
\item
  \textbf{Contested.} A weak link from a thought that says it replaces
  something marks the target \passthrough{\lstinline!contested!}. That
  covers a \passthrough{\lstinline!Correction!} or
  \passthrough{\lstinline!AssumptionInvalidated!} that refers to any
  thought, a \passthrough{\lstinline!Decision!} or
  \passthrough{\lstinline!StrategyShift!} that refers to an earlier
  decision, constraint, preference, plan or strategy with a plain
  reference or none, and any \passthrough{\lstinline!Contradicts!} edge.
  A contested thought stays visible, names the contesting thought, and
  that thought is promoted directly above it. Weak links do not hide,
  because a plain reference is often background and hiding it would lose
  valid memory.
\item
  \textbf{Annotated everywhere.} Every thought returned on every surface
  carries \passthrough{\lstinline!version\_status!}
  (\passthrough{\lstinline!current!},
  \passthrough{\lstinline!contested!} or
  \passthrough{\lstinline!superseded!}) and, when not current, its
  \passthrough{\lstinline!successors!} and
  \passthrough{\lstinline!current\_version!}, including results fetched
  with \passthrough{\lstinline!include\_invalidated=true!}.
\item
  \textbf{Advice at write time.} When a new
  \passthrough{\lstinline!Correction!},
  \passthrough{\lstinline!AssumptionInvalidated!} or
  \passthrough{\lstinline!Decision!} looks like a replacement but
  carries no strong relation, the append response returns
  \passthrough{\lstinline!warnings!} naming up to three likely targets,
  so the agent can add the typed link.
\end{itemize}

\passthrough{\lstinline!with\_resolve\_successors(false)!} on a library
query restores plain hiding.

A query with \passthrough{\lstinline!as\_of = T!} sees the chain as it
was at time \passthrough{\lstinline!T!}. It considers only thoughts
committed at or before \passthrough{\lstinline!T!}, and hides a thought
only if an invalidating edge was in force at
\passthrough{\lstinline!T!}: the edge's source was committed by then,
its \passthrough{\lstinline!valid\_at!} is not after
\passthrough{\lstinline!T!}, and its
\passthrough{\lstinline!invalid\_at!} is not at or before
\passthrough{\lstinline!T!}.

The live view treats every invalidating edge as permanent. An edge whose
\passthrough{\lstinline!invalid\_at!} has passed still hides its target
from live reads. Only \passthrough{\lstinline!as\_of!} reads honour edge
validity.

\subsubsection{3.7 Chains, Branches and Merges}

A chain is one independent log, named by a chain key. Keys are validated
at every entry point: not blank, at most 256 bytes, no control
characters (\passthrough{\lstinline!validate\_chain\_key!}). The chain
registry (\passthrough{\lstinline!mentisdb-registry.json!}) lists chains
with their schema version, storage adapter, thought and agent counts,
and location.

\passthrough{\lstinline!branch\_from!} creates a new chain whose first
thought is a \passthrough{\lstinline!StateSnapshot!} with a cross-chain
\passthrough{\lstinline!BranchesFrom!} edge to a thought of the source
chain. Nothing is copied and the source is untouched; the branch
inherits context by reference. Ranked search through the daemon follows
\passthrough{\lstinline!BranchesFrom!} edges and also searches the
ancestor chains, deduplicating by thought id. It searches each ancestor
whole, so thoughts appended to the parent after the fork point appear in
the branch's results too.

Merging is a server operation. It re-appends the source chain's
thoughts, unsigned, onto the target and then deregisters the source. It
copies type, role, content, importance, confidence, tags and concepts,
and maps agents by id. It drops relations, refs, sessions, entity types,
source episodes, signatures and original timestamps, and it refuses
targets that are not open-write chains. A merge therefore preserves what
was remembered but not how it was connected or who signed it.


\subsection{4. Integrity and Access Control}

\subsubsection{4.1 The Hash Chain}

Each thought's \passthrough{\lstinline!hash!} is the lowercase hex
SHA-256 of a canonical encoding of the record
(\passthrough{\lstinline!compute\_thought\_hash!}). The canonical form
is a fixed struct serialized with bincode 2 in its standard
configuration. It contains every field except
\passthrough{\lstinline!hash!} itself and two fields excluded for
backward compatibility: \passthrough{\lstinline!entity\_type!} and
\passthrough{\lstinline!source\_episode!}. It includes
\passthrough{\lstinline!prev\_hash!} and the signature fields. Every
thought therefore satisfies:

\begin{lstlisting}
thought.hash      == sha256(bincode(canonical(thought)))
thought.prev_hash == previous_thought.hash      ("" for index 0)
thought.index     == its position
\end{lstlisting}

Every open verifies the whole chain
(\passthrough{\lstinline!verify\_integrity!}): it checks each index,
each \passthrough{\lstinline!prev\_hash!} link, and recomputes each
hash. This costs O(n) and runs before the chain serves reads. Checking
links alone would miss an edited record, which is why the hashes are
recomputed.

This is tamper \emph{evidence}. An attacker who can write the file can
rewrite it entirely and recompute every hash. Evidence becomes useful
when the head hash is recorded somewhere the attacker cannot reach, and
signatures (Section 4.2) bind individual records to keys the attacker
does not hold.

Two gaps are known. First, \passthrough{\lstinline!entity\_type!} and
\passthrough{\lstinline!source\_episode!} are not covered by the record
hash, so they can be edited in the file without breaking the chain (v2
signatures do cover them). Second, migrating a chain from schema V0 to
V2 recomputes all hashes without checking the old ones, so tamper
evidence does not extend across that migration.

\subsubsection{4.2 Signatures}

A thought may carry a detached Ed25519 signature (RFC 8032, strict
verification) and the id of the signing key. The signature is checked on
append against an active, unrevoked key of the thought's agent, on every
chain. A bad signature is always refused, whatever the chain's policy.

Two payload versions exist, distinguished by the stored signature
itself:

{\def\LTcaptype{none} % do not increment counter
\begin{longtable}[]{@{}
  >{\raggedright\arraybackslash}p{(\linewidth - 4\tabcolsep) * \real{0.2766}}
  >{\raggedright\arraybackslash}p{(\linewidth - 4\tabcolsep) * \real{0.0851}}
  >{\raggedright\arraybackslash}p{(\linewidth - 4\tabcolsep) * \real{0.6383}}@{}}
\toprule\noalign{}
\begin{minipage}[b]{\linewidth}\raggedright
Stored \passthrough{\lstinline!thought\_signature!}
\end{minipage} & \begin{minipage}[b]{\linewidth}\raggedright
Version
\end{minipage} & \begin{minipage}[b]{\linewidth}\raggedright
Payload
\end{minipage} \\
\midrule\noalign{}
\endhead
\bottomrule\noalign{}
\endlastfoot
64 bytes & v1 & Legacy JSON of the semantic fields. Not bound to a chain
or position, so replayable; accepted only on open chains. \\
65 bytes starting \passthrough{\lstinline!0x02!} & v2 & Domain
\passthrough{\lstinline!mentisdb/thought/v2!}, then chain key, agent,
\passthrough{\lstinline!prev\_hash!}, session, type, role, content,
confidence, importance, tags (which include the scope), concepts, refs,
relations with validity, entity type and source episode. \\
\end{longtable}
}

The v2 payload (\passthrough{\lstinline!encode\_thought\_v2!} in
\passthrough{\lstinline!src/signing.rs!}) uses big-endian
length-prefixed fields and a NUL-terminated domain tag, so a signature
for one payload type can never verify as another. Binding the chain key
and the previous head means a signed thought is valid at exactly one
position of one chain. The cost is a race: if another writer moves the
head first, the append fails with the typed error
\passthrough{\lstinline!StaleHead!} and the client signs again against
the new head. The \passthrough{\lstinline!id!},
\passthrough{\lstinline!index!} and \passthrough{\lstinline!timestamp!}
are assigned by the chain and covered by the record hash rather than the
signature.

Signing happens on the client. The CLI signs locally with
\passthrough{\lstinline!--private-key-file!} and never sends the seed.

\subsubsection{4.3 Write Policies Recorded in the Chain}

Each chain has a write policy (\passthrough{\lstinline!AccessMode!} in
\passthrough{\lstinline!src/access.rs!}):

{\def\LTcaptype{none} % do not increment counter
\begin{longtable}[]{@{}
  >{\raggedright\arraybackslash}p{(\linewidth - 2\tabcolsep) * \real{0.2105}}
  >{\raggedright\arraybackslash}p{(\linewidth - 2\tabcolsep) * \real{0.7895}}@{}}
\toprule\noalign{}
\begin{minipage}[b]{\linewidth}\raggedright
Mode
\end{minipage} & \begin{minipage}[b]{\linewidth}\raggedright
A new thought must
\end{minipage} \\
\midrule\noalign{}
\endhead
\bottomrule\noalign{}
\endlastfoot
\passthrough{\lstinline!Open!} (default) & nothing more; unsigned writes
are accepted \\
\passthrough{\lstinline!Signed!} & carry a valid v2 signature by an
active key registered to its agent \\
\passthrough{\lstinline!Allowlist!} & carry a valid v2 signature by a
key in the policy's allowed set that is still active in the registry \\
\end{longtable}
}

The policy is not configuration. It lives in the chain as signed
\passthrough{\lstinline!AccessPolicy!} thoughts, replayed at open. Each
record holds the mode, the governor key set, an RFC 6962 Merkle root
over that set, a sequence number and the hash of the previous policy
record. A change after the first record must be signed by a current
governor. The first record is trust-on-first-use: it is accepted only
from an operator-authenticated entry point (the library, a global bearer
token, or the dashboard PIN), and its signer becomes the first governor.
An invalid record found on replay is ignored for enforcement and
reported by \passthrough{\lstinline!mentisdb chain access verify!}.

\passthrough{\lstinline!Signed!} chains trust the agent registry's key
list, which is a sidecar outside the hash chain.
\passthrough{\lstinline!Allowlist!} chains do not, because their allowed
keys are committed in the chain. Policies govern writes only; they do
not restrict reads. The full design, formats and threat model are in
\passthrough{\lstinline!docs/CHAIN\_ACCESS.md!}.

\subsubsection{4.4 Transport Authentication}

The daemon binds to \passthrough{\lstinline!127.0.0.1!} by default.
Bearer tokens are optional and off by default
(\passthrough{\lstinline!src/auth.rs!}). When enabled, tokens are stored
as SHA-256 hashes in a \passthrough{\lstinline!0600!} file and compared
in constant time. A token is either global or scoped to one or more
chains. REST and MCP share one table that marks each operation as
chain-scoped or global-only, normalize chain keys before authorization,
and authorize the same fields the handler executes. Chain-scoped tools
that name no chain require access to the default chain.

While tokens are off, requests that come from a browser (they carry
\passthrough{\lstinline!Origin!} or
\passthrough{\lstinline!Sec-Fetch-Site!}) must name an allowed host,
which blocks DNS-rebinding and cross-site requests to the loopback
listeners. The dashboard sits behind an optional PIN with a per-IP limit
of five failures per five minutes, sends a strict
Content-Security-Policy and frame-denial headers, and guards
state-changing calls against cross-site requests.

Several authorization bypasses have been found and fixed, in 0.10.9 and
again in the unreleased changes, including decoy chain-key fields that
let a chain-scoped token read another chain and a fail-open default for
unlisted MCP tools. Deployments that expose the daemon beyond loopback
should run current master and enable tokens.


\subsection{5. Storage and Durability}

\subsubsection{5.1 Record Format}

Each chain is one append-only file,
\passthrough{\lstinline!<stem>.tcbin!}. Each record is an 8-byte
little-endian length followed by the bincode 2 encoding of the
\passthrough{\lstinline!Thought!}
(\passthrough{\lstinline!persist\_binary\_thought!}). A record payload
may not exceed 10 MiB.

The file stem is the chain key with characters outside
\passthrough{\lstinline![A-Za-z0-9\_-]!} replaced by
\passthrough{\lstinline!\_!}, a dash, and the first 16 hex characters of
the SHA-256 of the chain key. Two keys that sanitize to the same string
therefore still get different files.

A legacy JSON-lines format (\passthrough{\lstinline!.jsonl!}) can still
be read for migration but no new chain uses it.

\subsubsection{5.2 File Layout}

\begin{lstlisting}
$MENTISDB_DIR/                         (default ~/.cloudllm/mentisdb)
  mentisdb-registry.json               chain registry
  mentisdb-skills.bin                  skill registry
  mentisdb-webhooks.json               webhook registrations
  bearer-tokens.json                   hashed bearer tokens (0600)
  tls/cert.pem, tls/key.pem            self-signed certificate
  .fastembed_cache/                    local embedding model
  <stem>.tcbin                         the hash chain (canonical)
  <stem>.agents.json                   agent registry
  <stem>.entity-types.json             entity-type registry
  <stem>.vectors.<model>.<ver>.<dim>d.json       vector sidecar snapshot
  <stem>.vectors.<model>.<ver>.<dim>d.json.wal   its write-ahead log
  <stem>.vectors.<model>.<ver>.<dim>d.hnsw.bin   persisted HNSW graph
  <stem>.vectors.managed.json          which sidecars stay in sync on append
  <stem>.auto_edges.bin                implicit-edge overlay
  migrations/vX_to_vY/                 archived pre-migration files
\end{lstlisting}

\subsubsection{5.3 Write Path and Durability Modes}

Appends go through a dedicated writer thread per chain. Two modes exist:

{\def\LTcaptype{none} % do not increment counter
\begin{longtable}[]{@{}
  >{\raggedright\arraybackslash}p{(\linewidth - 4\tabcolsep) * \real{0.2547}}
  >{\raggedright\arraybackslash}p{(\linewidth - 4\tabcolsep) * \real{0.3727}}
  >{\raggedright\arraybackslash}p{(\linewidth - 4\tabcolsep) * \real{0.3727}}@{}}
\toprule\noalign{}
\begin{minipage}[b]{\linewidth}\raggedright
Mode
\end{minipage} & \begin{minipage}[b]{\linewidth}\raggedright
Behaviour
\end{minipage} & \begin{minipage}[b]{\linewidth}\raggedright
What a crash can lose
\end{minipage} \\
\midrule\noalign{}
\endhead
\bottomrule\noalign{}
\endlastfoot
Strict (\passthrough{\lstinline!auto\_flush = true!}, the default) &
Each append returns after its group commit is written to the operating
system. When other appends are already queued, the writer waits up to 2
ms (\passthrough{\lstinline!MENTISDB\_GROUP\_COMMIT\_MS!}) to batch
them. & Nothing on a process crash. On power loss, writes since the last
sync unless \passthrough{\lstinline!MENTISDB\_FSYNC=1!}. \\
Buffered (\passthrough{\lstinline!auto\_flush = false!}) & Appends are
queued (capacity 128) and written every 16 records. & Queued and
unwritten records. \\
\end{longtable}
}

\passthrough{\lstinline!MENTISDB\_FSYNC=1!} adds one
\passthrough{\lstinline!fsync!} per group commit, and registry files are
then synced before their atomic rename. Explicit flush, shutdown and
drop always sync. A failed group commit is rolled back to the pre-commit
file length, so a full disk cannot leave a torn record followed by its
duplicate.

One writer process per data directory is assumed. There is no
cross-process file lock.

\subsubsection{5.4 Crash Recovery}

Opening a chain repairs only what is safe to repair.

\begin{itemize}
\tightlist
\item
  \textbf{Torn tail.} If the final record is strictly partial (a short
  length prefix, or a payload that ends at end of file), it is copied to
  \passthrough{\lstinline!<file>.torn-<millis>!} with mode
  \passthrough{\lstinline!0600!} and truncated before the writer reopens
  the file.
\item
  \textbf{Mid-file corruption.} If the bytes after a bad length prefix
  still decode as a complete record with a 64-digit hex hash, the damage
  is not a torn tail. The open fails with
  \passthrough{\lstinline!InvalidData!} and the file is left untouched,
  so no acknowledged thought is ever discarded as ``tail''.
\item
  \textbf{Bounded decoding.} Every decode of untrusted bytes runs under
  a budget: 64 MiB per thought record and 1 GiB per derived file
  (\passthrough{\lstinline!src/decode\_limits.rs!}). A forged length
  prefix yields \passthrough{\lstinline!InvalidData!} instead of an
  allocation that aborts the process.
\end{itemize}

These rules are pinned by adversarial tests that cut the last record at
every byte
(\passthrough{\lstinline!tests/core\_crash\_recovery\_tests.rs!}).

\subsubsection{5.5 Derived State}

Derived files make search fast, and each one can be rebuilt from the
log.

\begin{itemize}
\tightlist
\item
  \textbf{Vector sidecars} store one embedding per thought for one
  model, version and dimension. Changing any of these selects a
  different file, so incompatible embeddings never mix. Each append
  writes one integrity-chained record to the sidecar's write-ahead log
  (magic \passthrough{\lstinline!MDBVWAL1!}). The snapshot is compacted
  when the log reaches one sixteenth of the entries, and at least 32. A
  sidecar is fresh only if its thought count, head hash and entry count
  all match the chain.
\item
  \textbf{Background rebuild.} Opening a chain never embeds its history
  on the request path. A missing or stale sidecar is queued to one
  process-wide worker, one job per sidecar at a time, built in batches
  of 32 with one shared model, and adopted when done. Until then ranked
  search serves lexical and graph signals. Chains of 32 thoughts or
  fewer embed inline.
\item
  \textbf{Implicit-edge overlay} (Section 6.6) records the thought count
  and head hash it covers. On load it replays a missing tail of up to
  2048 thoughts, and otherwise rebuilds.
\item
  \textbf{HNSW graphs} are reloaded only when their stored count matches
  the chain.
\end{itemize}

\subsubsection{5.6 Schema Versions and Migration}

{\def\LTcaptype{none} % do not increment counter
\begin{longtable}[]{@{}
  >{\raggedright\arraybackslash}p{(\linewidth - 2\tabcolsep) * \real{0.5000}}
  >{\raggedright\arraybackslash}p{(\linewidth - 2\tabcolsep) * \real{0.5000}}@{}}
\toprule\noalign{}
\begin{minipage}[b]{\linewidth}\raggedright
Version
\end{minipage} & \begin{minipage}[b]{\linewidth}\raggedright
Change
\end{minipage} \\
\midrule\noalign{}
\endhead
\bottomrule\noalign{}
\endlastfoot
V0 & Original layout, no version field (JSON lines) \\
V1 & Explicit version, registry-backed chains, optional signatures \\
V2 & \passthrough{\lstinline!Reframe!} type,
\passthrough{\lstinline!Supersedes!} relation, cross-chain
\passthrough{\lstinline!chain\_key!} on relations \\
V3 & \passthrough{\lstinline!valid\_at!} and
\passthrough{\lstinline!invalid\_at!} on relations; later
\passthrough{\lstinline!entity\_type!} and
\passthrough{\lstinline!source\_episode!} were appended without a
version bump \\
\end{longtable}
}

The loader tries the current layout first and falls back record by
record to the legacy layouts, so a file with mixed versions still loads.
When the daemon starts it migrates registered chains before serving: it
writes the new file beside the old one, archives the old file under
\passthrough{\lstinline!migrations/!}, and renames the new file into
place. A chain written with the older JSON-based hash is rehashed with
the bincode hash on open. Both steps skip work that is already done, so
running them twice is harmless. Both also block: a large chain is
unavailable until its migration finishes.

Because bincode encodes enum variants by position, new variants must be
appended at the end of an enum and existing ones never reordered. A
chain that contains a variant added later, such as an
\passthrough{\lstinline!AccessPolicy!} record, cannot be read by an
older MentisDB.

\subsubsection{5.7 Backups}

A backup is a ZIP file named
\passthrough{\lstinline!mentisdb-YYYY-MM-DD-HH-MM-SS.mentis!} holding a
manifest with the SHA-256 and size of each file, the registries, chain
files and sidecars, skills, webhooks and optionally the TLS material.
Restore checks the manifest before writing, bounds sizes and file
counts, rejects symlinks and zip bombs, and never splices raw records
into a chain with a \passthrough{\lstinline!Signed!} or
\passthrough{\lstinline!Allowlist!} policy; such a backup is imported as
a separate, renamed chain.


\subsection{6. Retrieval}

Two read paths exist. \emph{Filter search}
(\passthrough{\lstinline!query!}) returns thoughts that match indexed
fields and a case-insensitive substring, in append order. \emph{Ranked
search} (\passthrough{\lstinline!query\_ranked!}) scores candidates with
several signals and returns them best first. Agents use ranked search
almost always. This section describes ranked search with its defaults as
of current master.

\subsubsection{6.1 Candidates}

Ranking starts from a candidate set
(\passthrough{\lstinline!retrieval\_candidates!}). It removes
invalidated thoughts unless
\passthrough{\lstinline!include\_invalidated!} is set, applies
\passthrough{\lstinline!as\_of!} (Section 3.6), and applies filters on
type, role, agent, tags, concepts, scope, entity type, session and time
range. Filters only restrict; they never add score. After scoring,
superseded thoughts that would have ranked are replaced by their live
successors and contested thoughts are preceded by their contesters
(Section 3.6).

\subsubsection{6.2 Lexical Signal: BM25 with Document-Frequency Gating}

Text is normalized by lowercasing, tokenizing and stemming with the
Snowball English (Porter2) stemmer. At query time, a table of 325
irregular forms maps words such as ``went'' and ``saw'' to their lemmas
(``go'', ``see''), which stemming cannot do.

Each thought is indexed in five fields, scored with BM25
(\passthrough{\lstinline!k1 = 1.2!}, \passthrough{\lstinline!b = 0.75!})
and weighted:

{\def\LTcaptype{none} % do not increment counter
\begin{longtable}[]{@{}
  >{\raggedright\arraybackslash}p{(\linewidth - 4\tabcolsep) * \real{0.7639}}
  >{\raggedright\arraybackslash}p{(\linewidth - 4\tabcolsep) * \real{0.1111}}
  >{\raggedright\arraybackslash}p{(\linewidth - 4\tabcolsep) * \real{0.1250}}@{}}
\toprule\noalign{}
\begin{minipage}[b]{\linewidth}\raggedright
Field
\end{minipage} & \begin{minipage}[b]{\linewidth}\raggedright
Weight
\end{minipage} & \begin{minipage}[b]{\linewidth}\raggedright
DF cutoff
\end{minipage} \\
\midrule\noalign{}
\endhead
\bottomrule\noalign{}
\endlastfoot
content & 1.0 & 30\% \\
tags & 1.6 & 30\% \\
concepts & 1.4 & 30\% \\
agent id & 1.5 & 70\% \\
agent registry text (name, owner, description, aliases) & 1.1 & 60\% \\
\end{longtable}
}

\begin{lstlisting}
idf(t)        = ln( (N - df + 0.5) / (df + 0.5) + 1 )
score(t, f)   = idf(t) * tf*(k1+1) / (tf + k1*(1 - b + b*len_f/avg_len_f)) * weight_f
\end{lstlisting}

MentisDB has no stop-word list. Instead, for each query term and each
field, the term is skipped in that field when
\passthrough{\lstinline!df / N!} exceeds the field's cutoff. This
\emph{document-frequency gate} is a hard cliff rather than a
down-weight: a term in 29\% of thoughts scores fully in
\passthrough{\lstinline!content!}, and one in 31\% scores nothing there.
A gated term can still score through a field whose cutoff it respects.
Details that matter in practice:

\begin{itemize}
\tightlist
\item
  By default \passthrough{\lstinline!df!} counts thoughts that contain
  the term in \emph{any} field (global scope), and that one ratio is
  compared with each field's cutoff. Registry text is attached to every
  thought of an agent, so a word in an agent's description counts as
  present in all of that agent's thoughts. On one real 750-thought
  chain, project words that occur in 2\% to 9\% of the content were
  gated out of it by registry text alone.
\item
  A per-field scope that counts each field separately exists as an
  opt-in (\passthrough{\lstinline!MENTISDB\_DF\_SCOPE=field!}, or
  \passthrough{\lstinline!df\_scope!} on the library query). It is not
  the default because it has not yet been run through LoCoMo and
  LongMemEval.
\item
  Gating is off for chains under 20 thoughts.
  \passthrough{\lstinline!N!} counts every committed thought, including
  invalidated ones and ones excluded by filters. The IDF of a surviving
  term always uses the global \passthrough{\lstinline!df!}.
\end{itemize}

\textbf{Synonyms.} The daemon expands query terms with a built-in
thesaurus of 895 headwords, keyed by stem so inflected forms expand too.
Each term gets at most three synonyms at weight 0.3
(\passthrough{\lstinline!apply\_thesaurus\_if\_text!} in
\passthrough{\lstinline!src/server.rs!}). The library does not expand by
default; embedders pass synonyms explicitly.

\subsubsection{6.3 Vector Signal and Fusion}

The default build includes local embeddings through fastembed: the
\passthrough{\lstinline!all-MiniLM-L6-v2!} model, 384 dimensions, run on
ONNX on the CPU. A dependency-free \passthrough{\lstinline!local-text!}
provider (256 dimensions) is the fallback.
\passthrough{\lstinline!MENTISDB\_EMBEDDINGS!} selects
\passthrough{\lstinline!auto!} (default),
\passthrough{\lstinline!fastembed!},
\passthrough{\lstinline!local-text!} or \passthrough{\lstinline!off!}.

The query is embedded and compared by cosine with each thought's vector.
Up to 256 hits per sidecar are kept, and hits under cosine 0.04 are
dropped. The vector score is \passthrough{\lstinline!0.5 * cosine!}. It
is fused with the lexical score by:

\begin{lstlisting}
vector_contribution = vector * (1 + 35 * exp(-lexical / 3))
\end{lstlisting}

With no lexical match the vector signal is amplified 36 times, so a
paraphrase with no shared words can still rank. As the lexical score
grows the amplification decays smoothly, so vectors help moderate
lexical matches without overriding strong ones.

Below 50,000 vectors the search is an exact scan. At or above that
threshold (\passthrough{\lstinline!MENTISDB\_HNSW\_THRESHOLD!}) it
switches to an HNSW graph (M 48, ef\_construction 400, ef\_search 128).
A filtered HNSW search scans exactly when the allowed set is small (at
most 4096 thoughts or under 5\% of the corpus) and otherwise
oversamples, so filters do not lose matches the graph skipped.

\subsubsection{6.4 Graph Expansion}

Graph expansion is off unless the query sets
\passthrough{\lstinline!graph!}. When on, the top 20 lexical hits seed a
breadth-first traversal over explicit relations and the implicit-edge
overlay, bidirectional by default, to depth 2, visiting at most 128
thoughts (\passthrough{\lstinline!expand\_with\_overlay\_classified!}).
Each reached thought gets three scores:

\begin{itemize}
\tightlist
\item
  \passthrough{\lstinline!graph = 1 / depth!},
\item
  \passthrough{\lstinline!seed\_support = 0.2 * min(paths - 1, 5)!} for
  thoughts reached from several seeds,
\item
  \passthrough{\lstinline!relation = boost(kind) / depth!} for the
  strongest edge kind on its path.
\end{itemize}

{\def\LTcaptype{none} % do not increment counter
\begin{longtable}[]{@{}ll@{}}
\toprule\noalign{}
Relation kind & Boost \\
\midrule\noalign{}
\endhead
\bottomrule\noalign{}
\endlastfoot
ContinuesFrom & 0.60 \\
BranchesFrom & 0.55 \\
Corrects, Invalidates & 0.50 \\
Supersedes & 0.45 \\
DerivedFrom & 0.40 \\
Summarizes, CausedBy & 0.20 \\
Supports, Contradicts & 0.15 \\
RelatedTo (implicit) & 0.08 \\
References & 0.06 \\
\end{longtable}
}

\subsubsection{6.5 Session Cohesion, Importance, Confidence and Recency}

\textbf{Session cohesion} surfaces the turns around a good match, which
often hold the answer without sharing its words. Every candidate with a
lexical score of at least 3.0 is a seed. A candidate whose own lexical
score is under 5.0 and that lies within 12 positions of a seed receives
\passthrough{\lstinline!1.2 * (1 - d / 12)!}, where
\passthrough{\lstinline!d!} is the distance in append order. This is
always on (\passthrough{\lstinline!apply\_session\_cohesion!}).

\textbf{Importance} shifts close races towards what the author marked as
important: \passthrough{\lstinline!lexical * (importance - 0.5) * 0.3!}
when there is a lexical match, else
\passthrough{\lstinline!(importance - 0.5) * 0.1!}. \textbf{Confidence}
adds \passthrough{\lstinline!0.1 * confidence!}. \textbf{Recency} adds
at most 0.05, linear in append position; there is no time-based decay.
Recency is deliberately a tie-breaker so that old memories are not
pushed down.

The total is:

\begin{lstlisting}
total = lexical + vector_contribution + graph + relation + seed_support
      + importance + confidence + recency + session_cohesion (+ rrf)
\end{lstlisting}

\subsubsection{6.6 Implicit Edges}

Most agents write few explicit relations, which leaves graph expansion
little to walk. The implicit-edge overlay adds
\passthrough{\lstinline!RelatedTo!} edges between thoughts whose
embeddings have cosine at least 0.85, keeping the 5 strongest per
thought (\passthrough{\lstinline!MENTISDB\_AUTO\_EDGE\_THRESHOLD!},
\passthrough{\lstinline!MENTISDB\_AUTO\_EDGE\_K!}). Up to 128 thoughts
the build compares all pairs. Above that it runs one nearest-neighbour
search per thought through the vector backend. Each append adds the new
thought's edges with one linear scan. Chains above 50,000 thoughts get
no overlay. The overlay is saved every 32 appends and is never part of
the hash chain.

\subsubsection{6.7 Optional Reciprocal Rank Fusion}

With \passthrough{\lstinline!enable\_reranking!}, the top 50 candidates
are ranked three times: by lexical score, by vector score, and by
graph-related score. Only candidates with a positive score enter each
list. Each candidate then gains a consensus bonus:

\begin{lstlisting}
rrf = 50 * sum over lists containing it of 1 / (60 + rank)
\end{lstlisting}

A candidate ranked first in all three lists gains about 2.5. RRF is off
by default. Before 0.10.9 it replaced the additive total and collapsed
rankings; it now only adds.

\subsubsection{6.8 Explaining a Ranking}

Each hit carries its full score breakdown:
\passthrough{\lstinline!lexical!}, \passthrough{\lstinline!vector!},
\passthrough{\lstinline!graph!}, \passthrough{\lstinline!relation!},
\passthrough{\lstinline!seed\_support!},
\passthrough{\lstinline!importance!},
\passthrough{\lstinline!confidence!}, \passthrough{\lstinline!recency!},
\passthrough{\lstinline!session\_cohesion!},
\passthrough{\lstinline!rrf!} and \passthrough{\lstinline!total!}, plus
the matched terms, the fields they matched in, and for graph hits the
distance, seed path count, relation kinds and path. A human or an agent
can see why any thought ranked where it did.

\emph{Context bundles}
(\passthrough{\lstinline!query\_context\_bundles!}) group results by
seed: each lexical seed comes with the thoughts graph expansion reached
from it and the path that reached them, so an agent can read supporting
memories in provenance order.

\subsubsection{6.9 Deduplication on Write}

Automatic deduplication is off unless
\passthrough{\lstinline!MENTISDB\_DEDUP\_THRESHOLD!} is set; 0.85 is a
suggested value, not a default. When on, each unsigned append compares
the new thought's normalized token set with the last 64 thoughts
(\passthrough{\lstinline!MENTISDB\_DEDUP\_SCAN\_WINDOW!}) by Jaccard
similarity. If the best match reaches the threshold, the new thought
gets a \passthrough{\lstinline!Supersedes!} edge to it, which hides the
older one by default. A prior thought the author rated more important
than the new one is never superseded, so a throwaway restatement cannot
hide a critical constraint. Signed thoughts are never deduplicated.

The window is a real limit: a fact re-recorded hundreds of thoughts
later is not caught. Section 9.4 shows that duplicates are a measured
problem on real chains.

\subsubsection{6.10 Standing Context and Recent Context}

Memory has two halves. The episodic half is what happened, in order. The
semantic half is what stays true. Recent context
(\passthrough{\lstinline!mentisdb\_recent\_context!}, 12 thoughts by
default) serves the first. It forgets: a preference from three months
ago has scrolled out of any practical window.

Standing context (\passthrough{\lstinline!MentisDb::standing\_context!},
\passthrough{\lstinline!src/standing.rs!}) serves the second. It selects
live thoughts of the standing types, ranks them by importance and then
recency, collapses restatements, and caps the result:

{\def\LTcaptype{none} % do not increment counter
\begin{longtable}[]{@{}
  >{\raggedright\arraybackslash}p{(\linewidth - 4\tabcolsep) * \real{0.4196}}
  >{\raggedright\arraybackslash}p{(\linewidth - 4\tabcolsep) * \real{0.4196}}
  >{\raggedright\arraybackslash}p{(\linewidth - 4\tabcolsep) * \real{0.1608}}@{}}
\toprule\noalign{}
\begin{minipage}[b]{\linewidth}\raggedright
Group
\end{minipage} & \begin{minipage}[b]{\linewidth}\raggedright
Types
\end{minipage} & \begin{minipage}[b]{\linewidth}\raggedright
Cap
\end{minipage} \\
\midrule\noalign{}
\endhead
\bottomrule\noalign{}
\endlastfoot
Constraint, Preference, User trait, Relationship, Goal, Decision, Lesson
& \passthrough{\lstinline!Constraint!},
\passthrough{\lstinline!PreferenceUpdate!},
\passthrough{\lstinline!UserTrait!},
\passthrough{\lstinline!RelationshipUpdate!},
\passthrough{\lstinline!Goal!}, \passthrough{\lstinline!Decision!},
\passthrough{\lstinline!LessonLearned!} & 8 per type, 40 in total \\
Handoff & \passthrough{\lstinline!Handoff!} & the newest one \\
Checkpoint & \passthrough{\lstinline!Checkpoint!} or
\passthrough{\lstinline!Summary!} type or role & the newest one \\
\end{longtable}
}

Only user-scoped or unscoped thoughts are standing. Within a group,
restatements with Jaccard similarity of 0.85 or more collapse to the
newest. A superseded constraint or decision is replaced by its live
successor, even when the successor is not itself a standing type such as
a \passthrough{\lstinline!Correction!}, and contested entries are
flagged in the prompt. The result is a compact prompt block delimited by
\passthrough{\lstinline!=== STANDING CONTEXT ===!} that the agent loads
at startup, next to the recent-context block. It is available as an MCP
tool, a REST route (\passthrough{\lstinline!POST /v1/standing-context!})
and a CLI command (\passthrough{\lstinline!mentisdb thought standing!}).

\subsubsection{6.11 What Is Not in the Default Path}

The crate contains components that ranked search does not use by
default, and this paper makes no claim based on them: personalized
PageRank (\passthrough{\lstinline!src/search/ppr.rs!}), pseudo-relevance
feedback (\passthrough{\lstinline!query\_expansion.rs!}), a query-intent
classifier (\passthrough{\lstinline!query\_intent.rs!}), embedding
synonyms, and an \passthrough{\lstinline!llm\_rerank!} flag that nothing
reads. LLM memory extraction (\passthrough{\lstinline!src/llm.rs!})
exists as an explicit, opt-in operation that writes
\passthrough{\lstinline!LLMExtracted!} thoughts.


\subsection{7. The Agent Protocol}

A memory store is only as good as the habits of the agents that use it.
MentisDB ships its operating protocol as a skill
(\passthrough{\lstinline!MENTISDB\_SKILL.md!}), served to MCP clients as
the resource \passthrough{\lstinline!mentisdb://skill/core!} and pointed
to by the server's \passthrough{\lstinline!initialize!} instructions.
Its main rules are these.

\textbf{Startup.} List chains and pick the project's chain, bootstrap it
(which writes a genesis checkpoint only on an empty chain), load the
skill, list agents and reuse an existing identity, load standing
context, load recent context, then write a checkpoint. Agents never
invent a new agent id; they ask the user.

\textbf{Search before write.} Before every append the agent runs a
ranked search for its proposed memory and makes three decisions. Should
it write at all, or is this already known? Should it correct, invalidate
or supersede an earlier thought instead of writing a parallel one? Which
one to three earlier thoughts should it link to, with which typed edge?
This keeps chains small and connected, and it is the main defence
against duplicates. Automatic deduplication is only the downstream half.
A change of policy or decision must carry a
\passthrough{\lstinline!Supersedes!}, \passthrough{\lstinline!Corrects!}
or \passthrough{\lstinline!Invalidates!} relation to the thought it
replaces. A plain reference only marks the old version contested, and if
the append response carries warnings the agent adds the typed link.

\textbf{Write triggers.} Write immediately on a lesson learned, a
decision, a constraint, a correction or invalidated assumption, a
checkpoint, a completed task, an open question or a hypothesis. One
strong memory is better than many weak ones.

\textbf{Recall triggers.} Search before answering when a question
touches the user's preferences, a past decision, a constraint or how
something was done before, or when the user says ``again'', ``remember''
or ``last time''. The agent checks each hit's
\passthrough{\lstinline!version\_status!} and answers from the current
version when a hit is contested or superseded. Searching with
\passthrough{\lstinline!include\_invalidated=true!} is for history, not
for finding the current answer. Never say ``I don't know your
preference'' without having searched.

Standing context covers what must be known at the start. Recall triggers
cover what becomes relevant later. Together they address the complaint
that users have to repeat themselves.


\subsection{8. Interfaces}

\textbf{MCP.} The daemon serves the Model Context Protocol over
streamable HTTP. A stdio mode lets clients that only speak stdio connect
through a local proxy that starts the daemon if needed. The server
exposes 48 tools covering bootstrap, append and retrospectives, filter
and ranked search, standing and recent context, context bundles,
traversal, federated search across chains, branches and merges, agents
and keys, access policies and signature verification, entity types,
import and export, vectors, webhooks and the skill registry.

\textbf{REST.} A versioned API under \passthrough{\lstinline!/v1/!}
mirrors the tools with 51 routes and can be embedded in another server
with \passthrough{\lstinline!rest\_router!}.

\textbf{CLI.} The \passthrough{\lstinline!mentisdb!} binary is also a
client for every MCP tool, REST route and dashboard action, with typed
argument validation and \passthrough{\lstinline!--help!} examples on
each command (\passthrough{\lstinline!docs/CLI\_PARITY.md!} maps each
one and a test keeps the map in sync). For example:

\begin{lstlisting}[language=bash]
mentisdb thought append "Use Postgres for the billing service" \
  --chain myproj --agent dev --type Decision
mentisdb search "which database did we pick for billing" --chain myproj
mentisdb thought standing --chain myproj
mentisdb chain access verify --chain myproj
\end{lstlisting}

\textbf{Dashboard.} A local HTTPS web interface browses chains,
thoughts, agents, skills and access policies, and manages settings
behind its PIN.

\textbf{Skill registry.} Agent instruction bundles are versioned in
\passthrough{\lstinline!mentisdb-skills.bin!}. The first version is
stored whole and later versions as diffs. A skill can be deprecated or
revoked with its history kept, or deleted permanently. An agent with
registered keys must sign its uploads. The registry is integrity-checked
as a file; it is not a hash chain.

\textbf{Webhooks and setup.} Webhooks notify external services of
appends. \passthrough{\lstinline!mentisdb setup!} configures MCP
clients: Claude Code, Claude Desktop, Codex, Copilot CLI, VS Code
Copilot, Gemini, Grok, OpenCode and Qwen.


\subsection{9. Evaluation}

\subsubsection{9.1 Setup}

Two public benchmarks measure retrieval recall: for each question, does
the gold evidence turn appear in the top \emph{k} thoughts? Both
harnesses are in the repository
(\passthrough{\lstinline!locomo-benches/!},
\passthrough{\lstinline!lme-benches/!}) and run against a live daemon
with a fresh chain.

\begin{itemize}
\tightlist
\item
  \textbf{LoCoMo-10P}: ten long multi-session conversations, 1977
  questions, split into single-hop (1554) and multi-hop (423). Turns are
  ingested with \passthrough{\lstinline!ContinuesFrom!} edges between
  consecutive turns, speaker and turn tags, and importance 0.8 or 0.2 by
  speaker. Queries use ranked search with graph expansion to depth 3. A
  hit is substring containment of the gold evidence in a returned
  thought.
\item
  \textbf{LongMemEval}: 500 questions over a chain of about 10,900
  turns, scored at R@5, R@10 and R@20.
\end{itemize}

Both use the default embedding model
(\passthrough{\lstinline!fastembed!} MiniLM, 384 dimensions) and the
daemon's automatic thesaurus.

\subsubsection{9.2 Results}

Release 0.10.9.54 (October 2026) is the latest full run:

{\def\LTcaptype{none} % do not increment counter
\begin{longtable}[]{@{}lllll@{}}
\toprule\noalign{}
Benchmark & Pipeline & R@5 & R@10 & R@20 \\
\midrule\noalign{}
\endhead
\bottomrule\noalign{}
\endlastfoot
LoCoMo-10P & default & 63.9\% & 73.5\% & \\
LoCoMo-10P & with RRF & 66.0\% & 74.1\% & \\
LongMemEval & default & 68.0\% & 72.8\% & 78.8\% \\
LongMemEval & with RRF & 68.6\% & & \\
\end{longtable}
}

Query throughput was about 384 queries per second on both. The changes
after 0.10.9.54 up to the date of this paper do not touch ranking
defaults, except where Section 6 says so.

Earlier full runs show where the gains came from:

{\def\LTcaptype{none} % do not increment counter
\begin{longtable}[]{@{}
  >{\raggedright\arraybackslash}p{(\linewidth - 6\tabcolsep) * \real{0.1034}}
  >{\raggedright\arraybackslash}p{(\linewidth - 6\tabcolsep) * \real{0.5517}}
  >{\raggedright\arraybackslash}p{(\linewidth - 6\tabcolsep) * \real{0.1724}}
  >{\raggedright\arraybackslash}p{(\linewidth - 6\tabcolsep) * \real{0.1724}}@{}}
\toprule\noalign{}
\begin{minipage}[b]{\linewidth}\raggedright
Release
\end{minipage} & \begin{minipage}[b]{\linewidth}\raggedright
Main change
\end{minipage} & \begin{minipage}[b]{\linewidth}\raggedright
LoCoMo-10P R@10
\end{minipage} & \begin{minipage}[b]{\linewidth}\raggedright
LongMemEval R@5
\end{minipage} \\
\midrule\noalign{}
\endhead
\bottomrule\noalign{}
\endlastfoot
0.8.0 & baseline & & 57.2\% \\
0.8.0 & Porter stemming & & 61.6\% \\
0.8.1 & smooth fusion, cohesion, DF gate & 74.2\% & 67.6\% \\
0.8.9 & irregular lemmas & 71.9\% & 66.8\% \\
0.9.9 & automatic thesaurus, real embeddings & 72.6\% & 66.8\% \\
0.10.9.54 & stemmed thesaurus keys, damped synonyms, RRF fix & 73.5\% &
68.0\% \\
\end{longtable}
}

On a two-conversation LoCoMo subset, 0.8.1 reached 88.7\% R@10. That
number is for the subset only and should not be compared with
full-benchmark results.

\subsubsection{9.3 Where Misses Come From}

A near-miss analysis of 0.8.5 on LoCoMo-10P bucketed the misses at R@10.
A quarter of them (25.8\%) were found by R@20, which is a ranking error.
Of the rest, 43.3\% were not found even by R@50, because the evidence
shares almost no words with the question. A similar analysis of 0.9.8 on
LongMemEval found that 65.7\% of misses were not in the top 20.
Multi-hop questions remain the weakest category on LoCoMo, about 17 to
20 points below single-hop. The lesson is that lexical retrieval has a
hard ceiling on paraphrase, and a small embedding model only partly
lifts it.

\subsubsection{9.4 Recall on Real Project Chains}

Benchmarks use conversations, not the memories coding agents actually
write. The \passthrough{\lstinline!examples/recall\_audit.rs!} tool
audits a copy of a real chain. It builds synthetic queries from the
chain's own thoughts, runs them through the production ranked search
with the daemon's defaults, and reports where each target lands,
bucketed by its age. Query kinds range from easy
(\passthrough{\lstinline!drop50!}: the target's words with half removed)
to hard (\passthrough{\lstinline!ref!}: the opening of a later thought
that refers to the target, and \passthrough{\lstinline!sem!}: the
opening of the target's nearest semantic neighbour among thoughts
sharing few of its words).

The audit of real chains found three things:

\begin{enumerate}
\def\labelenumi{\arabic{enumi}.}
\tightlist
\item
  \textbf{No age effect.} Recall did not fall with the age of the
  target. Older memories are found as reliably as recent ones once the
  query reaches them. Ranking does not bury the past; startup loading
  did, which standing context addresses.
\item
  \textbf{Duplicates crowd results.} Near-duplicate thoughts, written
  because agents did not search first or because a fact was restated
  outside the dedup window, took several of the top ten slots and pushed
  other relevant memories out.
\item
  \textbf{Paraphrase is weak.} Paraphrase-style queries recalled the
  target in the top ten only 43\% to 55\% of the time, against much
  higher recall for verbatim-style queries.
\end{enumerate}

These findings set the priorities in Section 10.

\subsubsection{9.5 Cost}

{\def\LTcaptype{none} % do not increment counter
\begin{longtable}[]{@{}
  >{\raggedright\arraybackslash}p{(\linewidth - 2\tabcolsep) * \real{0.6962}}
  >{\raggedright\arraybackslash}p{(\linewidth - 2\tabcolsep) * \real{0.3038}}@{}}
\toprule\noalign{}
\begin{minipage}[b]{\linewidth}\raggedright
Operation
\end{minipage} & \begin{minipage}[b]{\linewidth}\raggedright
Measurement
\end{minipage} \\
\midrule\noalign{}
\endhead
\bottomrule\noalign{}
\endlastfoot
Unsigned append, open chain (buffered, CPU only) & about 55 µs \\
Signed append (one Ed25519 verification) & about 75 to 80 µs \\
Append with a 3000-thought managed sidecar & 2.1 ms \\
Rare one-term lexical query, 100,000 documents & 1.8 µs \\
Three-term lexical query, 100,000 documents & 0.93 ms \\
100 vector queries, 10,000 x 128 dimensions, exact scan & 33.8 ms
(recall@10 100\%) \\
100 vector queries, 10,000 x 128 dimensions, HNSW & 8.5 ms (recall@10
91.4\%) \\
\end{longtable}
}

The first five figures are from the Criterion benches
\passthrough{\lstinline!access\_policy!} and
\passthrough{\lstinline!storage\_paths!} on Apple silicon in release
mode, as recorded in the changelog and
\passthrough{\lstinline!docs/CHAIN\_ACCESS.md!}. The vector figures are
from \passthrough{\lstinline!hnsw\_comparison!}, re-run for this paper
on an Apple M4 Max with synthetic unit vectors. HNSW is about four times
faster here; it only becomes the default backend at 50,000 vectors,
where an exact scan grows linearly. With
\passthrough{\lstinline!MENTISDB\_FSYNC=1!} the disk flush dominates
append cost.


\subsection{10. Limitations and Future Work}

\subsubsection{10.1 Limitations}

\begin{itemize}
\tightlist
\item
  \textbf{Paraphrase recall.} On real chains, paraphrased queries find
  the target in the top ten only about half the time. The default
  encoder is a small model whose cosine scores are compressed, and its
  training window is shorter than long thoughts.
\item
  \textbf{Duplicates.} Deduplication is off by default and sees only the
  last 64 thoughts. Restated facts accumulate and crowd results. The
  protocol's search-before-write rule is the main defence, and it
  depends on agent discipline.
\item
  \textbf{Global document frequency.} The default gate counts registry
  text as present in every thought of an agent, which removes useful
  project words from content scoring. The per-field alternative is not
  yet benchmarked.
\item
  \textbf{Currency depends on links.} Only typed
  \passthrough{\lstinline!Supersedes!},
  \passthrough{\lstinline!Corrects!} and
  \passthrough{\lstinline!Invalidates!} edges hide an old version.
  Untyped replacements are detected by a type-based rule and only marked
  contested, and a replacement with no link at all is caught only by the
  append-time warning.
\item
  \textbf{Live view ignores edge expiry.} An invalidating edge hides its
  target forever in live reads, even after its
  \passthrough{\lstinline!invalid\_at!}.
\item
  \textbf{Integrity gaps.} Two fields are outside the record hash,
  legacy migration rehashes without checking, and the agent registry is
  outside the chain, so \passthrough{\lstinline!Signed!} chains trust a
  sidecar. Integrity is local: there is no consensus and no external
  anchoring of head hashes.
\item
  \textbf{Reads are not access-controlled.} Write policies do not
  restrict who can read a chain. Bearer tokens are the only read control
  and are off by default.
\item
  \textbf{Blocking migrations.} Schema and hash migrations run before a
  chain serves and rewrite the whole file.
\item
  \textbf{Merges lose structure.} A merge drops relations, signatures
  and timestamps.
\item
  \textbf{Single writer.} Two processes writing one data directory are
  not detected.
\item
  \textbf{Benchmarks measure retrieval, not answers.} Recall at \emph{k}
  says whether the evidence was retrieved, not whether the agent used it
  well.
\end{itemize}

\subsubsection{10.2 Future Work}

\begin{itemize}
\tightlist
\item
  \textbf{Duplicate collapse in results:} fold near-duplicate clusters
  into one hit so that restatements stop crowding the top ten.
\item
  \textbf{Measuring valid-version retrieval} on benchmarks with changing
  facts, and on real chains with the correction audit
  (\passthrough{\lstinline!examples/correction\_audit.rs!}).
\item
  \textbf{Better paraphrase recall:} retune vector fusion, evaluate
  larger local encoders, and embed long thoughts in chunks.
\item
  \textbf{Per-field document frequency} as the default once benchmarked.
\item
  \textbf{Read access control:} private chains whose reads require
  authentication, with mode changes recorded in the chain, and key
  administration restricted to the operator.
\item
  \textbf{Non-blocking, recoverable migrations} and hash coverage of
  every field.
\item
  \textbf{Anchoring head hashes} externally so that a full rewrite of a
  chain becomes detectable.
\item
  \textbf{Answer-level evaluation} with an LLM reader, alongside recall.
\end{itemize}


\subsection{11. Related Work}

Several systems give agents long-term memory. Mem0 extracts and updates
memories with an LLM over a vector store. Graphiti, behind Zep, builds a
temporal knowledge graph with LLM extraction and validity intervals.
Letta, formerly MemGPT, lets the model page memory in and out of its own
context. These systems put an LLM on the write path, which gives them
structured facts at the cost of latency, money and a second model's
errors.

MentisDB makes the opposite trade. It keeps what the agent wrote,
verbatim and typed, adds no model to the write path, and puts the work
into retrieval, the write protocol and integrity. Its distinguishing
properties are the hash chain, signatures bound to chain position, write
policies governed inside the chain, standing context, and fully
decomposed scores, all in one embedded binary. Validity intervals on
relations are shared with Graphiti's design. Hybrid lexical and vector
retrieval is common to most systems.


\subsection{12. Conclusion}

MentisDB treats agent memory as a verifiable log of typed thoughts and
builds retrieval, currency and governance on top of it without an
external service. Its defaults are stated in this paper exactly, its
benchmark numbers come with their setup, and its weaknesses are
measured: paraphrase recall and duplicates, not age, are what limit
recall on real chains. That is where the next work goes.


\subsection{References}

\begin{itemize}
\tightlist
\item
  Robertson, S. and Zaragoza, H. \emph{The Probabilistic Relevance
  Framework: BM25 and Beyond}. Foundations and Trends in Information
  Retrieval, 3(4), 2009.
\item
  Cormack, G. V., Clarke, C. L. A. and Büttcher, S. \emph{Reciprocal
  Rank Fusion outperforms Condorcet and individual Rank Learning
  Methods}. SIGIR, 2009.
\item
  Porter, M. F. \emph{An algorithm for suffix stripping}. Program,
  14(3), 1980.
\item
  Jaccard, P. \emph{Étude comparative de la distribution florale dans
  une portion des Alpes et des Jura}. Bulletin de la Société Vaudoise
  des Sciences Naturelles, 1901.
\item
  Malkov, Y. A. and Yashunin, D. A. \emph{Efficient and robust
  approximate nearest neighbor search using Hierarchical Navigable Small
  World graphs}. IEEE TPAMI, 2018.
\item
  Bernstein, D. J., Duif, N., Lange, T., Schwabe, P. and Yang, B.-Y.
  \emph{High-speed high-security signatures}. Journal of Cryptographic
  Engineering, 2012.
\item
  Josefsson, S. and Liusvaara, I. \emph{Edwards-Curve Digital Signature
  Algorithm (EdDSA)}. RFC 8032, 2017.
\item
  Laurie, B., Langley, A. and Kasper, E. \emph{Certificate
  Transparency}. RFC 6962, 2013.
\item
  FIPS PUB 180-4. \emph{Secure Hash Standard (SHS)}. NIST, 2015.
\item
  Reimers, N. and Gurevych, I. \emph{Sentence-BERT: Sentence Embeddings
  using Siamese BERT-Networks}. EMNLP, 2019.
\item
  Maharana, A. et al.~\emph{Evaluating Very Long-Term Conversational
  Memory of LLM Agents (LoCoMo)}. ACL, 2024.
\item
  Wu, D. et al.~\emph{LongMemEval: Benchmarking Chat Assistants on
  Long-Term Interactive Memory}. 2024.
\item
  Anthropic. \emph{Model Context Protocol (MCP) Specification}. 2024.
\end{itemize}

\end{document}
